\chapter{Transformations et changements d'échelle}
\label{exercices:16}
\section*{Exercice 16.1 -- Images de points}
Dans un repère orthonormé, $A(2;1)$. Donner son image par la symétrie d'axe horizontal, la symétrie de centre $O$, la translation de déplacement $(3;-2)$ et la rotation de centre $O$ de $90\degre$ antihoraire. Faire un schéma.

\section*{Exercice 16.2 -- Rapport négatif}
Une homothétie de rapport $-1{,}5$ transforme une figure de périmètre $8$ et d'aire $20$. Calculer le périmètre et l'aire images. Dans l'espace, quel serait le volume image d'un solide de volume $40$ ? Expliquer le signe négatif.

\section*{Exercice 16.3 -- Retrouver une échelle}
Deux figures semblables ont des aires $36$ et $81\unite{cm^2}$. Quel est le facteur de longueur de la première vers la seconde ? Que devient un côté de $4\unite{cm}$ ?

\section*{Exercice 16.4 -- Rosace et pavage}
Une rosace comporte $8$ motifs régulièrement répartis autour du même centre. Donner l'angle de rotation. Expliquer pourquoi six triangles équilatéraux de même côté peuvent se rejoindre autour d'un sommet sans trou. Un agrandissement de facteur $2$ conserve-t-il les aires ?
